% Seminar 2 - Pagina 7. Fragment pentru % Seminar 2 - Pagina 7. Fragment pentru % Seminar 2 - Pagina 7. Fragment pentru % Seminar 2 - Pagina 7. Fragment pentru \input{seminars/pages/seminar2-pagina7-ro}.
% Fara preambul si fara intrerupere fortata de pagina.
% NOTE EDITORIALE:
% - Primele randuri continua exercitiul 10 de pe pagina 6: surjectivitatea
%   corespondentei dintre submultimi si functiile caracteristice.
% - La exercitiul 9, textul manuscrisului argumenteaza prin multimea M
%   si noteaza apoi imposibilitatea unei injectii P(X) -> X; acest ultim
%   pas nu este demonstrat separat pe pagina.
% - Formula incluziunii-excluziunii si inductia au fost pastrate in
%   ordinea din manuscris. Ultimul calcul se opreste la marginea paginii.
% - Semnele taiate si schita din marginea opusa nu sunt reconstruite.

\subsection*{Seminar 2 - Pagina 7}

\noindent\underline{Sol. 10 (continuare).} $\varphi$ este surjectivă.

\noindent Fie $\varphi:X\longrightarrow\{0,1\}$. Caut $A\subseteq X$ a.î.
$\varphi=\varphi_A$. Aleg
\[
 A:=\varphi^{-1}(1)=\{x\in X\mid\varphi(x)=1\}.
\]
\noindent Atunci
\[
 \varphi_A=\varphi_{\varphi^{-1}(1)}=\varphi,
\]
\[
 \varphi_{\varphi^{-1}(1)}(y)=
 \begin{cases}
  1,& y\in\varphi^{-1}(1),\\
  0,& y\notin\varphi^{-1}(1)
 \end{cases}
 =
 \begin{cases}
  1,&\varphi(y)=1,\\
  0,&\varphi(y)=0.
 \end{cases}
\]

\bigskip
\noindent\fbox{9}\quad\underline{Sol.} Nu există
$f:X\longrightarrow\mathcal{P}(X)$ surjectivă.

\noindent Presupunem prin reducere la absurd că există
$f:X\longrightarrow\mathcal{P}(X)$, $f$ surjectivă.

\noindent Fie
\[
 M:=\{a\in X\mid a\notin f(a)\}=f(b),\qquad b\in X.
\]
\noindent Atunci
\[
 b\in M\ \Longleftrightarrow\ b\notin f(b)=M,\qquad\text{absurd},
\]
\[
 b\notin M\ \Longleftrightarrow\ b\in f(b)=M.
\]
\noindent Asta implică că nu există
$i:\mathcal{P}(X)\longrightarrow X$ injectivă.

\noindent $\mathcal{P}(X)$ nu se poate scufunda în $X$.

\bigskip
\noindent\underline{Principiul includerii şi excluderii.}

\noindent Fie $A_1,A_2,\ldots,A_n$ mulţimi finite. Atunci
\[
\begin{aligned}
 \bigl|A_1\cup\cdots\cup A_n\bigr|
 ={}&\sum_{i=1}^{n}|A_i|
 -\sum_{1\leq i<j\leq n}|A_i\cap A_j|\\
 &+\sum_{1\leq i<j<k\leq n}|A_i\cap A_j\cap A_k|
 -\cdots\\
 &+(-1)^{n-1}|A_1\cap\cdots\cap A_n|.
\end{aligned}
\]

\noindent\underline{Sol.} Inducţie după $n$.

\noindent $n=1$: O.K., $|A_1|=|A_1|$.

\noindent $n=2$:
\[
 |A_1\cup A_2|=|A_1|+|A_2|-|A_1\cap A_2|,
 \qquad\text{O.K.}
\]

\noindent $n\longmapsto n+1$:
\[
\begin{aligned}
 &|A_1\cup\cdots\cup A_n\cup A_{n+1}|\\
 &\quad=|A_1\cup\cdots\cup A_n|+|A_{n+1}|
 -|(A_1\cup\cdots\cup A_n)\cap A_{n+1}|\\
 &\quad=\sum_{i=1}^{n+1}|A_i|
 -\sum_{1\leq i<j\leq n}|A_i\cap A_j|
 +\sum_{1\leq i<j<k\leq n}|A_i\cap A_j\cap A_k|
 -\cdots
\end{aligned}
\]
% Calculul inductiv continua dincolo de sfarsitul paginii manuscrisului.
.
% Fara preambul si fara intrerupere fortata de pagina.
% NOTE EDITORIALE:
% - Primele randuri continua exercitiul 10 de pe pagina 6: surjectivitatea
%   corespondentei dintre submultimi si functiile caracteristice.
% - La exercitiul 9, textul manuscrisului argumenteaza prin multimea M
%   si noteaza apoi imposibilitatea unei injectii P(X) -> X; acest ultim
%   pas nu este demonstrat separat pe pagina.
% - Formula incluziunii-excluziunii si inductia au fost pastrate in
%   ordinea din manuscris. Ultimul calcul se opreste la marginea paginii.
% - Semnele taiate si schita din marginea opusa nu sunt reconstruite.

\subsection*{Seminar 2 - Pagina 7}

\noindent\underline{Sol. 10 (continuare).} $\varphi$ este surjectivă.

\noindent Fie $\varphi:X\longrightarrow\{0,1\}$. Caut $A\subseteq X$ a.î.
$\varphi=\varphi_A$. Aleg
\[
 A:=\varphi^{-1}(1)=\{x\in X\mid\varphi(x)=1\}.
\]
\noindent Atunci
\[
 \varphi_A=\varphi_{\varphi^{-1}(1)}=\varphi,
\]
\[
 \varphi_{\varphi^{-1}(1)}(y)=
 \begin{cases}
  1,& y\in\varphi^{-1}(1),\\
  0,& y\notin\varphi^{-1}(1)
 \end{cases}
 =
 \begin{cases}
  1,&\varphi(y)=1,\\
  0,&\varphi(y)=0.
 \end{cases}
\]

\bigskip
\noindent\fbox{9}\quad\underline{Sol.} Nu există
$f:X\longrightarrow\mathcal{P}(X)$ surjectivă.

\noindent Presupunem prin reducere la absurd că există
$f:X\longrightarrow\mathcal{P}(X)$, $f$ surjectivă.

\noindent Fie
\[
 M:=\{a\in X\mid a\notin f(a)\}=f(b),\qquad b\in X.
\]
\noindent Atunci
\[
 b\in M\ \Longleftrightarrow\ b\notin f(b)=M,\qquad\text{absurd},
\]
\[
 b\notin M\ \Longleftrightarrow\ b\in f(b)=M.
\]
\noindent Asta implică că nu există
$i:\mathcal{P}(X)\longrightarrow X$ injectivă.

\noindent $\mathcal{P}(X)$ nu se poate scufunda în $X$.

\bigskip
\noindent\underline{Principiul includerii şi excluderii.}

\noindent Fie $A_1,A_2,\ldots,A_n$ mulţimi finite. Atunci
\[
\begin{aligned}
 \bigl|A_1\cup\cdots\cup A_n\bigr|
 ={}&\sum_{i=1}^{n}|A_i|
 -\sum_{1\leq i<j\leq n}|A_i\cap A_j|\\
 &+\sum_{1\leq i<j<k\leq n}|A_i\cap A_j\cap A_k|
 -\cdots\\
 &+(-1)^{n-1}|A_1\cap\cdots\cap A_n|.
\end{aligned}
\]

\noindent\underline{Sol.} Inducţie după $n$.

\noindent $n=1$: O.K., $|A_1|=|A_1|$.

\noindent $n=2$:
\[
 |A_1\cup A_2|=|A_1|+|A_2|-|A_1\cap A_2|,
 \qquad\text{O.K.}
\]

\noindent $n\longmapsto n+1$:
\[
\begin{aligned}
 &|A_1\cup\cdots\cup A_n\cup A_{n+1}|\\
 &\quad=|A_1\cup\cdots\cup A_n|+|A_{n+1}|
 -|(A_1\cup\cdots\cup A_n)\cap A_{n+1}|\\
 &\quad=\sum_{i=1}^{n+1}|A_i|
 -\sum_{1\leq i<j\leq n}|A_i\cap A_j|
 +\sum_{1\leq i<j<k\leq n}|A_i\cap A_j\cap A_k|
 -\cdots
\end{aligned}
\]
% Calculul inductiv continua dincolo de sfarsitul paginii manuscrisului.
.
% Fara preambul si fara intrerupere fortata de pagina.
% NOTE EDITORIALE:
% - Primele randuri continua exercitiul 10 de pe pagina 6: surjectivitatea
%   corespondentei dintre submultimi si functiile caracteristice.
% - La exercitiul 9, textul manuscrisului argumenteaza prin multimea M
%   si noteaza apoi imposibilitatea unei injectii P(X) -> X; acest ultim
%   pas nu este demonstrat separat pe pagina.
% - Formula incluziunii-excluziunii si inductia au fost pastrate in
%   ordinea din manuscris. Ultimul calcul se opreste la marginea paginii.
% - Semnele taiate si schita din marginea opusa nu sunt reconstruite.

\subsection*{Seminar 2 - Pagina 7}

\noindent\underline{Sol. 10 (continuare).} $\varphi$ este surjectivă.

\noindent Fie $\varphi:X\longrightarrow\{0,1\}$. Caut $A\subseteq X$ a.î.
$\varphi=\varphi_A$. Aleg
\[
 A:=\varphi^{-1}(1)=\{x\in X\mid\varphi(x)=1\}.
\]
\noindent Atunci
\[
 \varphi_A=\varphi_{\varphi^{-1}(1)}=\varphi,
\]
\[
 \varphi_{\varphi^{-1}(1)}(y)=
 \begin{cases}
  1,& y\in\varphi^{-1}(1),\\
  0,& y\notin\varphi^{-1}(1)
 \end{cases}
 =
 \begin{cases}
  1,&\varphi(y)=1,\\
  0,&\varphi(y)=0.
 \end{cases}
\]

\bigskip
\noindent\fbox{9}\quad\underline{Sol.} Nu există
$f:X\longrightarrow\mathcal{P}(X)$ surjectivă.

\noindent Presupunem prin reducere la absurd că există
$f:X\longrightarrow\mathcal{P}(X)$, $f$ surjectivă.

\noindent Fie
\[
 M:=\{a\in X\mid a\notin f(a)\}=f(b),\qquad b\in X.
\]
\noindent Atunci
\[
 b\in M\ \Longleftrightarrow\ b\notin f(b)=M,\qquad\text{absurd},
\]
\[
 b\notin M\ \Longleftrightarrow\ b\in f(b)=M.
\]
\noindent Asta implică că nu există
$i:\mathcal{P}(X)\longrightarrow X$ injectivă.

\noindent $\mathcal{P}(X)$ nu se poate scufunda în $X$.

\bigskip
\noindent\underline{Principiul includerii şi excluderii.}

\noindent Fie $A_1,A_2,\ldots,A_n$ mulţimi finite. Atunci
\[
\begin{aligned}
 \bigl|A_1\cup\cdots\cup A_n\bigr|
 ={}&\sum_{i=1}^{n}|A_i|
 -\sum_{1\leq i<j\leq n}|A_i\cap A_j|\\
 &+\sum_{1\leq i<j<k\leq n}|A_i\cap A_j\cap A_k|
 -\cdots\\
 &+(-1)^{n-1}|A_1\cap\cdots\cap A_n|.
\end{aligned}
\]

\noindent\underline{Sol.} Inducţie după $n$.

\noindent $n=1$: O.K., $|A_1|=|A_1|$.

\noindent $n=2$:
\[
 |A_1\cup A_2|=|A_1|+|A_2|-|A_1\cap A_2|,
 \qquad\text{O.K.}
\]

\noindent $n\longmapsto n+1$:
\[
\begin{aligned}
 &|A_1\cup\cdots\cup A_n\cup A_{n+1}|\\
 &\quad=|A_1\cup\cdots\cup A_n|+|A_{n+1}|
 -|(A_1\cup\cdots\cup A_n)\cap A_{n+1}|\\
 &\quad=\sum_{i=1}^{n+1}|A_i|
 -\sum_{1\leq i<j\leq n}|A_i\cap A_j|
 +\sum_{1\leq i<j<k\leq n}|A_i\cap A_j\cap A_k|
 -\cdots
\end{aligned}
\]
% Calculul inductiv continua dincolo de sfarsitul paginii manuscrisului.
.
% Fara preambul si fara intrerupere fortata de pagina.
% NOTE EDITORIALE:
% - Primele randuri continua exercitiul 10 de pe pagina 6: surjectivitatea
%   corespondentei dintre submultimi si functiile caracteristice.
% - La exercitiul 9, textul manuscrisului argumenteaza prin multimea M
%   si noteaza apoi imposibilitatea unei injectii P(X) -> X; acest ultim
%   pas nu este demonstrat separat pe pagina.
% - Formula incluziunii-excluziunii si inductia au fost pastrate in
%   ordinea din manuscris. Ultimul calcul se opreste la marginea paginii.
% - Semnele taiate si schita din marginea opusa nu sunt reconstruite.

\subsection*{Seminar 2 - Pagina 7}

\noindent\underline{Sol. 10 (continuare).} $\varphi$ este surjectivă.

\noindent Fie $\varphi:X\longrightarrow\{0,1\}$. Caut $A\subseteq X$ a.î.
$\varphi=\varphi_A$. Aleg
\[
 A:=\varphi^{-1}(1)=\{x\in X\mid\varphi(x)=1\}.
\]
\noindent Atunci
\[
 \varphi_A=\varphi_{\varphi^{-1}(1)}=\varphi,
\]
\[
 \varphi_{\varphi^{-1}(1)}(y)=
 \begin{cases}
  1,& y\in\varphi^{-1}(1),\\
  0,& y\notin\varphi^{-1}(1)
 \end{cases}
 =
 \begin{cases}
  1,&\varphi(y)=1,\\
  0,&\varphi(y)=0.
 \end{cases}
\]

\bigskip
\noindent\fbox{9}\quad\underline{Sol.} Nu există
$f:X\longrightarrow\mathcal{P}(X)$ surjectivă.

\noindent Presupunem prin reducere la absurd că există
$f:X\longrightarrow\mathcal{P}(X)$, $f$ surjectivă.

\noindent Fie
\[
 M:=\{a\in X\mid a\notin f(a)\}=f(b),\qquad b\in X.
\]
\noindent Atunci
\[
 b\in M\ \Longleftrightarrow\ b\notin f(b)=M,\qquad\text{absurd},
\]
\[
 b\notin M\ \Longleftrightarrow\ b\in f(b)=M.
\]
\noindent Asta implică că nu există
$i:\mathcal{P}(X)\longrightarrow X$ injectivă.

\noindent $\mathcal{P}(X)$ nu se poate scufunda în $X$.

\bigskip
\noindent\underline{Principiul includerii şi excluderii.}

\noindent Fie $A_1,A_2,\ldots,A_n$ mulţimi finite. Atunci
\[
\begin{aligned}
 \bigl|A_1\cup\cdots\cup A_n\bigr|
 ={}&\sum_{i=1}^{n}|A_i|
 -\sum_{1\leq i<j\leq n}|A_i\cap A_j|\\
 &+\sum_{1\leq i<j<k\leq n}|A_i\cap A_j\cap A_k|
 -\cdots\\
 &+(-1)^{n-1}|A_1\cap\cdots\cap A_n|.
\end{aligned}
\]

\noindent\underline{Sol.} Inducţie după $n$.

\noindent $n=1$: O.K., $|A_1|=|A_1|$.

\noindent $n=2$:
\[
 |A_1\cup A_2|=|A_1|+|A_2|-|A_1\cap A_2|,
 \qquad\text{O.K.}
\]

\noindent $n\longmapsto n+1$:
\[
\begin{aligned}
 &|A_1\cup\cdots\cup A_n\cup A_{n+1}|\\
 &\quad=|A_1\cup\cdots\cup A_n|+|A_{n+1}|
 -|(A_1\cup\cdots\cup A_n)\cap A_{n+1}|\\
 &\quad=\sum_{i=1}^{n+1}|A_i|
 -\sum_{1\leq i<j\leq n}|A_i\cap A_j|
 +\sum_{1\leq i<j<k\leq n}|A_i\cap A_j\cap A_k|
 -\cdots
\end{aligned}
\]
% Calculul inductiv continua dincolo de sfarsitul paginii manuscrisului.
